\section{3.2 누적 보상, 할인율, 그리고 실습}
\begin{frame}{이 절의 흐름}
  3.1절에서 MDP의 다섯 요소 $(\mathcal{S}, \mathcal{A}, P, R,
  \gamma)$ 를 ``나열''하기만 했다. 이 절은 그중 \kb{$R$}(보상)과
  \kb{$\gamma$}(할인율)에 집중해서, 강화학습이 실제로 \textbf{최대화
  하는 대상 --- 리턴 --- 을 손으로$\cdot$코드로 계산할 수 있게 만든다.
  }\vskip0.6em
  수업의 흐름:
  \begin{enumerate}
    \item \kb{리턴}(할인된 누적 보상)의 정의와 유한/무한 경우의 계산법.
    \item ``할인율 0.95''가 실제로는 ``몇 스텝짜리 시야''를 뜻하는지
      --- \kb{유효 시계(effective horizon)}로 바꾸는 법.
    \item $\gamma$ 를 바꾸면 ``같은 환경에서 다른 행동을
      \textbf{최적으로}'' 만드는 예 --- 할인율이 단순 정규화
      장치가 아니라 \textbf{목표를 바꾸는} 매개변수임을 확인.
    \item Gymnasium에서 CartPole로 할인된 리턴을 실제로 누적.
  \end{enumerate}
\end{frame}
